\caption{\textbf{(a,b)} Three intervals of an illustrative simulated ARG yield matching local trees. All haplotypes supply alleles. References also supply ancestry vectors. Query ancestry is unknown. Ancestors carry age $\tau$ and sampled/root indicators. \textbf{(c)} $N_t$ counts nodes in tree $t$, $S$ sampled haplotypes and $d$ feature width. Two GNN passes each traverse every node upward then downward with shared weights. Generic forks summarize clade presence across the trees. Updates for gained, lost and consecutive clades use each node's own state. Tiny boxes show the actual 350-generation cut groups. The same five history rows persist, with $q_1$ in magenta. Open vectors show its current group context $f$. Filled rows show sample histories $z$, scanned independently in both directions with retention $\rho_t=\exp(-w_t/\lambda)$ and current weight $1-\rho_t$. Here $\lambda$ is a learned physical length scale. $\bar z$ averages both directions and the current group's sample rows. Subtracting $f$ and linearly projecting supplies a residual before pass two. Readout retains $S$ separate vectors $e_{tq}$ for each tree. \textbf{(d)} One shared BiGRU with one layer reads each sampled sequence with physical span $w_t$ and distances between neighboring centers $\Delta c_t^\pm$, transformed by $\log(1+\cdot)$. $\Vert$ concatenates local and recurrent embeddings. The MLP and softmax give ancestry probabilities $p_{tq}$, and argmax gives interval calls. Values and predictions are schematic.}
